%%%%%%%%%%%%%%%%%%%%%%%%%%%%%%%%%%%%%%%%%%%%%%%%%%%%%%%%%%%%%%%%%%%%%%%%%%%
% Función: organización mínima de las tesis por compendio de publicaciones.
%
% Uso: copia esta organización mínima a Book/content-compendium.tex.
% Modifica la copia para organizar el contenido de tu trabajo.
%%%%%%%%%%%%%%%%%%%%%%%%%%%%%%%%%%%%%%%%%%%%%%%%%%%%%%%%%%%%%%%%%%%%%%%%%%%

\part{\myCompendiumSummaryPartTitle}
\label{part:compendium-summary}

%%%%%%%%%%%%%%%%%%%%%%%%%%%%%%%%%%%%%%%%%%%%%%%%%%%%%%%%%%%%%%%%%%%%%%%%%%%
% Función: capítulo de ejemplo de introducción de la tesis por compendio.
%
% Uso: modifica este fichero para adaptar el contenido a tu trabajo.
%%%%%%%%%%%%%%%%%%%%%%%%%%%%%%%%%%%%%%%%%%%%%%%%%%%%%%%%%%%%%%%%%%%%%%%%%%%

\chapter{Introducción}
\label{cha:compendium-introduction}

Como ejemplo concreto de presentación de un área de investigación, puedes contrastar la localización visual en espacios inteligentes~\cite{Sanabria2026Landmark}, la evaluación funcional en actividades cotidianas~\cite{nieva2024proposal} y la localización de barcos mediante sensado acústico distribuido~\cite{ramirez2026jstars}. Son ejemplos reales de líneas de investigación distintas, no publicaciones atribuidas al doctorando ficticio de esta plantilla.

\section{Contexto y motivación}
\subsection{Contexto de la investigación}

Utiliza este capítulo para presentar el contexto de la investigación, su motivación y el problema central abordado por la tesis. Explica por qué el problema es relevante, delimita el alcance del trabajo y proporciona la información necesaria para comprender el resumen extendido. Este ejemplo también muestra el uso de acrónimos: \gls{AI} y \gls{HMI}.

\subsection{Planteamiento del problema}

Define el problema concreto abordado por la tesis, su alcance y las hipótesis que delimitan la investigación propuesta.

\section{Organización de la tesis}

Describe brevemente el contenido del resumen extendido y explica cómo las publicaciones incluidas en la parte~\ref{part:compendium-publications} sustentan el hilo argumental de la tesis.

%%% Local Variables:
%%% TeX-master: "../../book"
%%% End:

%%%%%%%%%%%%%%%%%%%%%%%%%%%%%%%%%%%%%%%%%%%%%%%%%%%%%%%%%%%%%%%%%%%%%%%%%%%
% Función: capítulo de ejemplo de objetivos de la tesis por compendio.
%
% Uso: modifica este fichero para adaptar el contenido a tu trabajo.
%%%%%%%%%%%%%%%%%%%%%%%%%%%%%%%%%%%%%%%%%%%%%%%%%%%%%%%%%%%%%%%%%%%%%%%%%%%

\chapter{Objetivos}
\label{cha:compendium-objectives}

Al formular objetivos medibles, distingue la tarea de estimación ilustrada por la localización de cara y boca~\cite{Sanabria2026Landmark}, el problema de evaluación ligado a la aplicación~\cite{nieva2024proposal} y el despliegue de sensores y análisis descritos para EYEFUL~\cite{nieva2023eyeful}. Úsalos como ejemplos de distintos niveles de objetivos, no como objetivos o resultados que tu tesis ya haya alcanzado.

\section{Objetivo general}

Expón la hipótesis principal o el objetivo general de la tesis y deriva los objetivos específicos abordados por las publicaciones. Haz explícita la relación entre estos objetivos y la organización del resumen extendido.

\section{Objetivos específicos}
\subsection{Objetivos científicos}

Identifica los avances de conocimiento perseguidos por la tesis y relaciona cada uno con las preguntas de investigación.

\subsection{Objetivos técnicos y de validación}

Describe los objetivos de implementación, experimentación y validación que permiten evaluar los objetivos científicos.

%%% Local Variables:
%%% TeX-master: "../../book"
%%% End:

%%%%%%%%%%%%%%%%%%%%%%%%%%%%%%%%%%%%%%%%%%%%%%%%%%%%%%%%%%%%%%%%%%%%%%%%%%%
% Función: capítulo de ejemplo de trabajos previos de la tesis por compendio.
%
% Uso: modifica este fichero para adaptar el contenido a tu trabajo.
%%%%%%%%%%%%%%%%%%%%%%%%%%%%%%%%%%%%%%%%%%%%%%%%%%%%%%%%%%%%%%%%%%%%%%%%%%%

\chapter{Trabajo previo}
\label{cha:compendium-previous-work}

Una revisión bibliográfica puede distinguir un método concreto de localización visual~\cite{Sanabria2026Landmark}, una propuesta de evaluación funcional~\cite{nieva2024proposal} y una aplicación de detección y localización mediante sensado acústico distribuido~\cite{ramirez2026jstars}. La revisión de métodos de aprendizaje automático para vigilancia de tuberías~\cite{tejedor2017machine} ofrece además un ejemplo de revisión del estado del arte, en lugar de un artículo sobre un método individual. Estos ejemplos ilustran formas de organizar las fuentes, pero no implican que todos los trabajos aborden la misma pregunta de investigación.

\section{Antecedentes científicos}
\subsection{Enfoques fundamentales}

Revisa aquí los antecedentes científicos y técnicos de la tesis. Organiza la discusión alrededor de las preguntas de investigación en lugar de presentar una lista aislada de publicaciones, e identifica las limitaciones que motivan el trabajo doctoral.

\subsection{Desarrollos recientes}

Resume los enfoques más relacionados con la tesis y compara sus hipótesis, fortalezas y limitaciones.

\section{Oportunidad de investigación}

Concluye la revisión identificando las cuestiones no resueltas que abordan los objetivos de esta tesis.

%%% Local Variables:
%%% TeX-master: "../../book"
%%% End:

%%%%%%%%%%%%%%%%%%%%%%%%%%%%%%%%%%%%%%%%%%%%%%%%%%%%%%%%%%%%%%%%%%%%%%%%%%%
% Función: capítulo de ejemplo de metodología de la tesis por compendio.
%
% Uso: modifica este fichero para adaptar el contenido a tu trabajo.
%%%%%%%%%%%%%%%%%%%%%%%%%%%%%%%%%%%%%%%%%%%%%%%%%%%%%%%%%%%%%%%%%%%%%%%%%%%

\chapter{Metodología}
\label{cha:compendium-methodology}

Para ilustrar la documentación de la metodología, fíjate en cómo un estudio de localización visual concreta su problema de estimación~\cite{Sanabria2026Landmark}, cómo EYEFUL describe el despliegue de sensores y el análisis de datos~\cite{nieva2023eyeful}, y cómo una aplicación de sensado acústico distribuido plantea la detección y la localización~\cite{ramirez2026jstars}. La publicación del conjunto de datos asociado~\cite{ramireztorres2026distributedacousticsensingdataset} es una fuente independiente para comentar la procedencia y reutilización de los datos. Estas referencias son ilustrativas. La tabla numérica que aparece a continuación es solo un ejemplo de formato y no reproduce sus resultados.

\section{Diseño de la investigación}
\subsection{Enfoque general}

Describe la metodología común que conecta las publicaciones, incluidos el diseño experimental, los datos, los procedimientos de análisis y los criterios de validación. Remite a las publicaciones incluidas para los detalles, manteniendo este capítulo como una explicación coherente de la metodología de la tesis.

\subsection{Datos y protocolo experimental}

Explica cómo se obtienen, preparan y dividen los datos, e identifica las variables y los controles utilizados en el conjunto de las publicaciones.

\section{Metodología de evaluación}

La figura~\ref{fig:compendium-example} y la tabla~\ref{tab:compendium-example} son ejemplos sencillos que puedes sustituir por material de tu tesis.

\begin{figure}[htbp]
  \centering
  \includegraphics[width=0.55\textwidth]{Figure1}
  \caption{Ejemplo de figura en el resumen extendido.}
  \label{fig:compendium-example}
\end{figure}

\begin{table}[htbp]
  \centering
  \renewcommand{\arraystretch}{1.2}
  \caption{Ejemplo de comparación de resultados experimentales.}
  \label{tab:compendium-example}
  \begin{tabular}{|l|c|c|}
    \hline
    Método &
    Métrica A &
    Métrica B \\
    \hline
    Referencia & 0.72 & 0.68 \\
    \hline
    Método propuesto & 0.84 & 0.79 \\
    \hline
  \end{tabular}
\end{table}

%%% Local Variables:
%%% TeX-master: "../../book"
%%% End:

%%%%%%%%%%%%%%%%%%%%%%%%%%%%%%%%%%%%%%%%%%%%%%%%%%%%%%%%%%%%%%%%%%%%%%%%%%%
% Función: capítulo de ejemplo de contribuciones de la tesis por compendio.
%
% Uso: modifica este fichero para adaptar el contenido a tu trabajo.
%%%%%%%%%%%%%%%%%%%%%%%%%%%%%%%%%%%%%%%%%%%%%%%%%%%%%%%%%%%%%%%%%%%%%%%%%%%

\chapter{Contribuciones de la tesis}
\label{cha:compendium-contributions}

Como ejemplos para identificar una contribución, puedes comparar el método de localización~\cite{Sanabria2026Landmark}, la propuesta de evaluación funcional~\cite{nieva2024proposal} y el estudio de detección y localización de barcos~\cite{ramirez2026jstars}. Cita el trabajo real cuando lo utilices como antecedente y distingue su autoría de tus propias contribuciones. Estas referencias reales son independientes de las tres publicaciones ficticias utilizadas a continuación para mostrar la estructura del compendio.

\section{Resumen de las contribuciones}

Presenta y discute las contribuciones originales de la tesis. Explica cómo se relacionan las publicaciones entre sí, qué objetivo aborda cada una y cómo sus resultados combinados constituyen una contribución doctoral unificada. Los tres artículos ficticios distribuidos con este ejemplo muestran cómo citar cada publicación incluida al discutir su contribución~\cite{garcia2026paper1,garcia2026paper2,garcia2026paper3}.

\subsection{Relación entre las publicaciones}

Explica el papel de cada publicación y la progresión desde la pregunta inicial de investigación hasta los resultados conjuntos de la tesis.

\section{Material de apoyo a la investigación}
\subsection{Vídeos y demostraciones}

Recuerda que puedes incluir muy fácilmente enlaces a vídeos en el documento y generar automáticamente un listado similar al de tablas, figuras o algoritmos. La orden se define en \texttt{Config/postamble.tex}, mientras que el listado se genera desde \texttt{cover/list-of-videos.tex}. Para usarla, basta con que incluyas elementos similares a \videoLink{Amplificadores realimentados: motivación}{https://youtu.be/9xXuasS8B3M} o \videoLink{Amplificadores realimentados: teoría ideal}{https://youtu.be/ZGMwGX7VCbg}.

\subsection{Algoritmos y código fuente}

Y también puedes incluir algoritmos como el algoritmo~\ref{alg:restriction}, y símbolos como~\ac{xidet}.

Además de listados de código fuente como en el listado~\ref{cod:sample2}.


\begin{algorithm}
  \caption{IntervalRestriction\label{IR}}
  \label{alg:restriction}
  \DontPrintSemicolon
  % \dontprintsemicolon
  \KwData{$G=(X,U)$ such that $G^{tc}$ is an order.}
  \KwResult{$G'=(X,V)$ with $V\subseteq U$ such that $G'^{tc}$ is an
    interval order.}
  \Begin{
    $V \longleftarrow U$\;
    $S \longleftarrow \emptyset$\;
    \For{$x\in X$}{
      $NbSuccInS(x) \longleftarrow 0$\;
      $NbPredInMin(x) \longleftarrow 0$\;
      $NbPredNotInMin(x) \longleftarrow |ImPred(x)|$\;
    }
    \For{$x \in X$}{
      \If{$NbPredInMin(x) = 0$ {\bf and} $NbPredNotInMin(x) = 0$}{
        $AppendToMin(x)$}
    }
    \nl\While{$S \neq \emptyset$}{\label{InRes1}
      \nlset{REM} remove $x$ from the list of $T$ of maximal index\;\label{InResR}
      \lnl{InRes2}\While{$|S \cap ImSucc(x)| \neq |S|$}{
        \For{$ y \in S-ImSucc(x)$}{
          \{ remove from $V$ all the arcs $zy$ : \}\;
          \For{$z \in ImPred(y) \cap Min$}{
            remove the arc $zy$ from $V$\;
            $NbSuccInS(z) \longleftarrow NbSuccInS(z) - 1$\;
            move $z$ in $T$ to the list preceding its present list\;
            \{i.e. If $z \in T[k]$, move $z$ from $T[k]$ to
            $T[k-1]$\}\;
          }
          $NbPredInMin(y) \longleftarrow 0$\;
          $NbPredNotInMin(y) \longleftarrow 0$\;
          $S \longleftarrow S - \{y\}$\;
          $AppendToMin(y)$\;
        }
      }
      $RemoveFromMin(x)$\;
    }
  }
\end{algorithm}


\begin{codefloat}
\lstinputlisting[style=Cnice]{appendix/hello.c}
\caption{Ejemplo de código fuente con estilo \texttt{Cnice}, de nuevo
  con un \texttt{lstinputlisting} dentro de un \texttt{codefloat}}
\label{cod:sample2}
\end{codefloat}

%%% Local Variables:
%%% TeX-master: "../../book"
%%% End:

%%%%%%%%%%%%%%%%%%%%%%%%%%%%%%%%%%%%%%%%%%%%%%%%%%%%%%%%%%%%%%%%%%%%%%%%%%%
% Función: capítulo de ejemplo de conclusiones de la tesis por compendio.
%
% Uso: modifica este fichero para adaptar el contenido a tu trabajo.
%%%%%%%%%%%%%%%%%%%%%%%%%%%%%%%%%%%%%%%%%%%%%%%%%%%%%%%%%%%%%%%%%%%%%%%%%%%

\chapter{Conclusiones y líneas futuras}
\label{cha:compendium-conclusions}

Como ejemplo de conexión entre conclusiones y alcance, compara los objetivos ligados a la aplicación en evaluación funcional~\cite{nieva2024proposal}, el contexto del sistema descrito para EYEFUL~\cite{nieva2023eyeful} y el escenario de sensado diferente de la localización de barcos~\cite{ramirez2026jstars}. Utiliza la comparación para preguntarte qué afirmaciones están sustentadas por tus propias evidencias y cuáles exigirían más validación. No presentes resultados de esas publicaciones como resultados de esta tesis ficticia.

\section{Conclusiones principales}

Resume las conclusiones sustentadas por el conjunto de la tesis, relaciónalas con los objetivos y discute las principales limitaciones. Cierra el resumen extendido identificando las líneas de investigación que se derivan de los resultados.

\section{Limitaciones}

Discute las hipótesis y restricciones prácticas que limitan la generalización de los resultados y deben considerarse al interpretar las conclusiones.

\section{Líneas futuras de investigación}
\subsection{Extensiones inmediatas}

Identifica las mejoras y experimentos que se derivan directamente del trabajo presentado en la tesis.

\subsection{Oportunidades de investigación a largo plazo}

Plantea cuestiones más amplias abiertas por el conjunto de las contribuciones y publicaciones.

%%% Local Variables:
%%% TeX-master: "../../book"
%%% End:


%%%%%%%%%%%%%%%%%%%%%%%%%%%%%%%%%%%%%%%%%%%%%%%%%%%%%%%%%%%%%%%%%%%%%%%%%%%
% Función: bibliografía.
%
% Uso: inclusión automática; normalmente no debes modificarlo.
% No se compila por separado.
%%%%%%%%%%%%%%%%%%%%%%%%%%%%%%%%%%%%%%%%%%%%%%%%%%%%%%%%%%%%%%%%%%%%%%%%%%%

% IMPORTANT: YOU DON'T HAVE TO EDIT THIS FILE, JUST EDIT THE bibliofiles.tex file
% IMPORTANT: YOU DON'T HAVE TO EDIT THIS FILE, JUST EDIT THE bibliofiles.tex file
% IMPORTANT: YOU DON'T HAVE TO EDIT THIS FILE, JUST EDIT THE bibliofiles.tex file
% IMPORTANT: YOU DON'T HAVE TO EDIT THIS FILE, JUST EDIT THE bibliofiles.tex file
% IMPORTANT: YOU DON'T HAVE TO EDIT THIS FILE, JUST EDIT THE bibliofiles.tex file
% IMPORTANT: YOU DON'T HAVE TO EDIT THIS FILE, JUST EDIT THE bibliofiles.tex file
% IMPORTANT: YOU DON'T HAVE TO EDIT THIS FILE, JUST EDIT THE bibliofiles.tex file
% IMPORTANT: YOU DON'T HAVE TO EDIT THIS FILE, JUST EDIT THE bibliofiles.tex file
% IMPORTANT: YOU DON'T HAVE TO EDIT THIS FILE, JUST EDIT THE bibliofiles.tex file
% IMPORTANT: YOU DON'T HAVE TO EDIT THIS FILE, JUST EDIT THE bibliofiles.tex file
% IMPORTANT: YOU DON'T HAVE TO EDIT THIS FILE, JUST EDIT THE bibliofiles.tex file
% IMPORTANT: YOU DON'T HAVE TO EDIT THIS FILE, JUST EDIT THE bibliofiles.tex file
% IMPORTANT: YOU DON'T HAVE TO EDIT THIS FILE, JUST EDIT THE bibliofiles.tex file
% IMPORTANT: YOU DON'T HAVE TO EDIT THIS FILE, JUST EDIT THE bibliofiles.tex file
% IMPORTANT: YOU DON'T HAVE TO EDIT THIS FILE, JUST EDIT THE bibliofiles.tex file
% IMPORTANT: YOU DON'T HAVE TO EDIT THIS FILE, JUST EDIT THE bibliofiles.tex file
% IMPORTANT: YOU DON'T HAVE TO EDIT THIS FILE, JUST EDIT THE bibliofiles.tex file
% IMPORTANT: YOU DON'T HAVE TO EDIT THIS FILE, JUST EDIT THE bibliofiles.tex file
% IMPORTANT: YOU DON'T HAVE TO EDIT THIS FILE, JUST EDIT THE bibliofiles.tex file
% IMPORTANT: YOU DON'T HAVE TO EDIT THIS FILE, JUST EDIT THE bibliofiles.tex file
% IMPORTANT: YOU DON'T HAVE TO EDIT THIS FILE, JUST EDIT THE bibliofiles.tex file


\ifthenelse{\equal{\bibliosystem}{biblatex}}
{
  % Use biblatex instead of bibtex

  %% Wen changing to biblatex, we could not define here the bibliography files, so that
  %% You need to edit the bibliofiles.tex file to do so

  %% % Now add all bib files to be processed
  %% \addbibresource{\myreferencespath\mybibfileOne}
  %% % \addbibresource{\myreferencespath\mybibfileTwo}
  %% % ...
  %% % \addbibresource{\myreferencespath\mybibfileN}

  % Other documents (such as the anteproyecto) retain their usual bibliography.
  \ifdefined\ThesisPrintFinalBibliography
    \ThesisPrintFinalBibliography
  \else
    \printbibliography[heading=bibintoc]
  \fi


}
{
  % Use bibtex

  %\bibliographystyle{plainnat}
  %\bibliographystyle{dinat}
  %\bibliographystyle{unsrt}
  \bibliographystyle{IEEEtran}

  % The following is overly complicated because I was not able to do so in
  % another way. The problem is the bibliography command being "called"
  % from both the root and anteproyecto directories...
  %
  %%%%%%%%%%%%%%%%%%%%%%%%%%%%%%%%%%%%%%%%%%%%%%%%%%%%%%%%%%%%%%%%%%%%%%%%%%%
% Función: selección de los ficheros de bibliografía del documento.
%
% Uso: modifica este fichero para incluir las definiciones de tu trabajo.
%%%%%%%%%%%%%%%%%%%%%%%%%%%%%%%%%%%%%%%%%%%%%%%%%%%%%%%%%%%%%%%%%%%%%%%%%%%

%% Here define as many bibfiles as needed
%%
%% It is compulsory that they are named as \mybibfileOne
%% \mybibfileTwo, \mybibfileThree, ... \mybibfileTen
%%
%% If you need more than ten, you will have to edit
%% Config/preamble.tex and Book/biblio/bibliography.tex
%% to support this adition
%%
%% The file names may change at your will, but they must
%% be in the Book/biblio directory

\newcommand{\mybibfileOne}{biblio/biblio.bib}
\newcommand{\mybibfileTwo}{biblio/nobiblio.bib}
%% \newcommand{\mybibfileThree}{AudioVisualNew.bib}
%% \newcommand{\mybibfileFour}{biblio/audiotracking.bib}
%% \newcommand{\mybibfileFive}{biblio/audiovisualtracking.bib}
%% \newcommand{\mybibfileSix}{biblio/backgroungsubstraction.bib}
%% \newcommand{\mybibfileSeven}{biblio/databases.bib}
%% \newcommand{\mybibfileEight}{biblio/evalmetrics.bib}
%% \newcommand{\mybibfileNine}{biblio/facedetect.bib}
%% \newcommand{\mybibfileTen}{biblio/facedetectADABOOST.bib}
%% \newcommand{\mybibfileEleven}{biblio/facedetectmultiview.bib}
%% \newcommand{\mybibfileTwelve}{biblio/facedetectprob2d.bib}
%% \newcommand{\mybibfileThirteen}{biblio/others.bib}
%% \newcommand{\mybibfileFourteen}{biblio/skindetect.bib}
%% \newcommand{\mybibfileFifteen}{biblio/tracking.bib}
%% \newcommand{\mybibfileSixteen}{biblio/videotracking.bib}
%% \newcommand{\mybibfileSeventeen}{biblio/voiceActivityDetection.bib}
%% \newcommand{\mybibfileEighteen}{biblio/headposeextraction.bib}
%% \newcommand{\mybibfileNineteen}{biblio/AudioVisualSpeakerTracking.bib}
%% \newcommand{\mybibfileTwenty}{biblio/BibliogPFVJ.bib}
%% \newcommand{\mybibfileTwentyone}{biblio/tools.bib}
%% \newcommand{\mybibfileTwentytwo}{biblio/infrared.bib}
%% \newcommand{\mybibfileTwentythree}{}
%% \newcommand{\mybibfileTwentyfour}{}
%% \newcommand{\mybibfileTwentyfive}{}


  \newcommand{\mybibfiles}{}
  \ifdef{\mybibfileOne}
  {
  \let\oldmybibfiles\mybibfiles
  \renewcommand{\mybibfiles}{\myreferencespath\mybibfileOne}
  }
  {
  \errorYOUmustDEFINEatLEASTmybibfileOneInbibliofilesDOTtex
  }
  \ifdef{\mybibfileTwo}
  {
  \let\oldmybibfiles\mybibfiles
  \renewcommand{\mybibfiles}{\oldmybibfiles,\myreferencespath\mybibfileTwo}
  }
  {
  }
  \ifdef{\mybibfileThree}
  {
  \let\oldmybibfiles\mybibfiles
  \renewcommand{\mybibfiles}{\oldmybibfiles,\myreferencespath\mybibfileThree}
  }
  {
  }

  \ifdef{\mybibfileFour}
  {
  \let\oldmybibfiles\mybibfiles
  \renewcommand{\mybibfiles}{\oldmybibfiles,\myreferencespath\mybibfileFour}
  }
  {
  }
  \ifdef{\mybibfileSix}
  {
  \let\oldmybibfiles\mybibfiles
  \renewcommand{\mybibfiles}{\oldmybibfiles,\myreferencespath\mybibfileSix}
  }
  {
  }

  \ifdef{\mybibfileSeven}
  {
  \let\oldmybibfiles\mybibfiles
  \renewcommand{\mybibfiles}{\oldmybibfiles,\myreferencespath\mybibfileSeven}
  }
  {
  }

  \ifdef{\mybibfileEight}
  {
  \let\oldmybibfiles\mybibfiles
  \renewcommand{\mybibfiles}{\oldmybibfiles,\myreferencespath\mybibfileEight}
  }
  {
  }

  \ifdef{\mybibfileNine}
  {
  \let\oldmybibfiles\mybibfiles
  \renewcommand{\mybibfiles}{\oldmybibfiles,\myreferencespath\mybibfileNine}
  }
  {
  }

  \ifdef{\mybibfileTen}
  {
  \let\oldmybibfiles\mybibfiles
  \renewcommand{\mybibfiles}{\oldmybibfiles,\myreferencespath\mybibfileTen}
  }
  {
  }

  \ifdef{\mybibfileEleven}
  {
    \let\oldmybibfiles\mybibfiles
    \renewcommand{\mybibfiles}{\oldmybibfiles,\myreferencespath\mybibfileEleven}
  }
  {
  }

  \ifdef{\mybibfileTwelve}
  {
    \let\oldmybibfiles\mybibfiles
    \renewcommand{\mybibfiles}{\oldmybibfiles,\myreferencespath\mybibfileTwelve}
  }
  {
  }

  \ifdef{\mybibfileThirteen}
  {
    \let\oldmybibfiles\mybibfiles
    \renewcommand{\mybibfiles}{\oldmybibfiles,\myreferencespath\mybibfileThirteen}
  }
  {
  }

  \ifdef{\mybibfileFourteen}
  {
    \let\oldmybibfiles\mybibfiles
    \renewcommand{\mybibfiles}{\oldmybibfiles,\myreferencespath\mybibfileFourteen}
  }
  {
  }

  \ifdef{\mybibfileFifteen}
  {
    \let\oldmybibfiles\mybibfiles
    \renewcommand{\mybibfiles}{\oldmybibfiles,\myreferencespath\mybibfileFifteen}
  }
  {
  }

  \ifdef{\mybibfileSixteen}
  {
    \let\oldmybibfiles\mybibfiles
    \renewcommand{\mybibfiles}{\oldmybibfiles,\myreferencespath\mybibfileSixteen}
  }
  {
  }

  \ifdef{\mybibfileSeventeen}
  {
    \let\oldmybibfiles\mybibfiles
    \renewcommand{\mybibfiles}{\oldmybibfiles,\myreferencespath\mybibfileSeventeen}
  }
  {
  }

  \ifdef{\mybibfileEighteen}
  {
    \let\oldmybibfiles\mybibfiles
    \renewcommand{\mybibfiles}{\oldmybibfiles,\myreferencespath\mybibfileEighteen}
  }
  {
  }

  \ifdef{\mybibfileNineteen}
  {
    \let\oldmybibfiles\mybibfiles
    \renewcommand{\mybibfiles}{\oldmybibfiles,\myreferencespath\mybibfileNineteen}
  }
  {
  }

  \ifdef{\mybibfileTwenty}
  {
    \let\oldmybibfiles\mybibfiles
    \renewcommand{\mybibfiles}{\oldmybibfiles,\myreferencespath\mybibfileTwenty}
  }
  {
  }

  \ifdef{\mybibfileTwentyone}
  {
    \let\oldmybibfiles\mybibfiles
    \renewcommand{\mybibfiles}{\oldmybibfiles,\myreferencespath\mybibfileTwentyone}
  }
  {
  }

  \ifdef{\mybibfileTwentytwo}
  {
    \let\oldmybibfiles\mybibfiles
    \renewcommand{\mybibfiles}{\oldmybibfiles,\myreferencespath\mybibfileTwentytwo}
  }
  {
  }

  \ifdef{\mybibfileTwentythree}
  {
    \let\oldmybibfiles\mybibfiles
    \renewcommand{\mybibfiles}{\oldmybibfiles,\myreferencespath\mybibfileTwentythree}
  }
  {
  }

  \ifdef{\mybibfileTwentyfour}
  {
    \let\oldmybibfiles\mybibfiles
    \renewcommand{\mybibfiles}{\oldmybibfiles,\myreferencespath\mybibfileTwentyfour}
  }
  {
  }

  \ifdef{\mybibfileTwentyfive}
  {
    \let\oldmybibfiles\mybibfiles
    \renewcommand{\mybibfiles}{\oldmybibfiles,\myreferencespath\mybibfileTwentyfive}
  }
  {
  }

  % Do not touch this
  % The commands around the \bibliography{} command should be included if using bibtex, but
  % according to Gonzalo Corral's PR on July 2022, they break compilation in overleaf. I think
  % this should not happen if the bib files are in iso-8859-1 when using bibtex instead of
  % biber, but need to be tested (TODO)
  \inputencoding{latin1}
  \bibliography{\mybibfiles}
  \inputencoding{utf8}
}

%%% Local Variables:
%%% TeX-master: "../book"
%%% coding: utf-8
%%% End:




\part{\myCompendiumPublicationsPartTitle}
\label{part:compendium-publications}

% Añade aquí una llamada a \includecompendiumpublication por cada publicación.
% El argumento entre corchetes es opcional y admite texto y órdenes de LaTeX.
% Los cinco argumentos entre llaves son obligatorios y aparecen en este orden:
% etiqueta interna, título abreviado para el índice y la cabecera, título
% completo, clave de biblio/biblio.bib y ruta del PDF de la publicación.
% Sustituye los datos del ejemplo y elimina los signos % para activarlo:
%
% \includecompendiumpublication[Estado de publicación, contribución del doctorando o permisos de reutilización]
%   {articulo-1}
%   {Título abreviado}
%   {Título completo de la publicación}
%   {clave-bibliografica}
%   {publications/articulo-1.pdf}

%%% Local Variables:
%%% TeX-master: "book"
%%% End:
